
\documentclass{article}
\usepackage{colortbl}
\usepackage{makecell}
\usepackage{multirow}
\usepackage{supertabular}

\begin{document}

\newcounter{utterance}

\centering \large Interaction Transcript for game `eqbench', experiment `eq\_bench', episode 148 with reflexion{-}Qwen3.5{-}9B.
\vspace{24pt}

{ \footnotesize  \setcounter{utterance}{1}
\setlength{\tabcolsep}{0pt}
\begin{supertabular}{c@{$\;$}|p{.15\linewidth}@{}p{.15\linewidth}p{.15\linewidth}p{.15\linewidth}p{.15\linewidth}p{.15\linewidth}}
    \# & \multicolumn{2}{c}{Player} && \multicolumn{2}{c}{Game Master} \\
    \hline

    \theutterance \stepcounter{utterance}  
    & & & \multicolumn{4}{p{0.6\linewidth}}{
        \cellcolor[rgb]{0.9,0.9,0.9}{
            \makecell[{{p{\linewidth}}}]{
                \texttt{\tiny{[P1$\langle$GM]}}
                \texttt{You have attempted this same task before and did not succeed. These are the plans you wrote after those attempts.} \\
\texttt{Trial \#1: New plan: I incorrectly identified "Disgusted" as the primary emotion because Jackson's final line ("Better than pretending to be something I'm not") sounded defensive and self{-}righteous. However, looking back at the dialogue, Jackson explicitly admitted, "You're right, I've never had to worry about surviving," which directly validates Alejandro's core argument about survival. This admission indicates Jackson finally comprehended Alejandro's perspective, making "Understood" the most accurate emotional response. Next time, I will prioritize identifying moments where a character explicitly concedes a point or validates the other's worldview as the strongest indicator of "Understood," rather than focusing on the tone of the final retort. I will re{-}evaluate the scores to reflect that Jackson feels understood (high score) and less disgusted (lower score) than previously thought.} \\
\\ 
\texttt{Trial \#2: New plan: I incorrectly scored "Understood" as 8 and "Disgusted" as 4, failing to recognize that Jackson's admission ("You're right, I've never had to worry about surviving") was a moment of genuine realization, not just a defensive retort. While Jackson feels understood, the intensity of his realization is the defining feature of the scene, making "Understood" the overwhelmingly dominant emotion compared to the others. Next time, I will assign a significantly higher score to "Understood" (e.g., 9 or 10) to reflect the explicit validation of Alejandro's worldview, and lower the scores for "Disgusted" and "Doubtful" to accurately represent that these are secondary or fading reactions in the face of this breakthrough moment. I will ensure the highest score clearly indicates the primary emotional shift.} \\
\\ 
\\ 
\\ 
\texttt{Your task is to predict the likely emotional responses of a character in this dialogue:} \\
\\ 
\texttt{Jackson: Sanchez. Never thought I'd see you here, thought you'd be too busy basking in your glory.} \\
\texttt{Alejandro: Riley, always the charmer. I could say the same about you. How's the knee?} \\
\texttt{Jackson: Better than your conscience, I reckon.} \\
\texttt{Alejandro: Ah, the ever{-}righteous Riley. You think you're the only one who's had to fight for everything?} \\
\texttt{Jackson: I don't recall stepping on others to get to the top.} \\
\texttt{Alejandro: And I don't recall asking for your opinion.} \\
\texttt{Jackson: You never did. Too busy making a spectacle of yourself.} \\
\texttt{Alejandro: And you're too busy playing the victim. } \\
\texttt{Jackson: I ain't the one parading around with another man's woman.} \\
\texttt{Alejandro: And I ain't the one who couldn't keep her.} \\
\texttt{Jackson: You think this is about her? This is about respect, something you wouldn't understand.} \\
\texttt{Alejandro: And you think this is about respect? This is about survival, something you've never had to worry about.} \\
\texttt{Jackson: You're right, I've never had to worry about surviving. I've been too busy living.} \\
\texttt{Alejandro: And how's that working out for you, Riley?} \\
\texttt{Jackson: Better than pretending to be something I'm not, Sanchez.} \\
\texttt{Alejandro: And who says I'm pretending? Maybe this is who I am.} \\
\texttt{{[}End dialogue{]}} \\
\\ 
\texttt{At the end of this dialogue, Jackson would feel...} \\
\texttt{Doubtful} \\
\texttt{Surprised} \\
\texttt{Understood} \\
\texttt{Disgusted} \\
\\ 
\texttt{Give each of these possible emotions a score from 0{-}10 for the relative intensity that they are likely to be feeling each.} \\
\\ 
\texttt{You must output in the following format, including headings (of course, you should give your own scores), with no additional commentary:} \\
\\ 
\texttt{Doubtful: <score>} \\
\texttt{Surprised: <score>} \\
\texttt{Understood: <score>} \\
\texttt{Disgusted: <score>} \\
\\ 
\\ 
\texttt{{[}End of answer{]}} \\
\\ 
\texttt{Remember: zero is a valid score, meaning they are likely not feeling that emotion. You must score at least one emotion > 0.} \\
\\ 
\texttt{Your answer:} \\
            }
        }
    }
     \\ \\

    \theutterance \stepcounter{utterance}  
    & \multicolumn{4}{p{0.6\linewidth}}{
        \cellcolor[rgb]{0.561,0.737,0.561}{
            \makecell[{{p{\linewidth}}}]{
                \texttt{\tiny{[P1$\rangle$GM]}}
                \texttt{Doubtful: 1} \\
\texttt{Surprised: 2} \\
\texttt{Understood: 9} \\
\texttt{Disgusted: 3} \\
            }
        }
    }
    & & \\ \\

    \theutterance \stepcounter{utterance}  
    & & & \multicolumn{2}{p{0.3\linewidth}}{
        \cellcolor[rgb]{0.9,0.9,0.9}{
            \makecell[{{p{\linewidth}}}]{
                \texttt{\tiny{[GM$|$GM]}}
                \texttt{Doubtful: 1} \\
\texttt{Surprised: 2} \\
\texttt{Understood: 9} \\
\texttt{Disgusted: 3} \\
            }
        }
    }
    & & \\ \\

    \theutterance \stepcounter{utterance}  
    & & & \multicolumn{2}{p{0.3\linewidth}}{
        \cellcolor[rgb]{0.9,0.9,0.9}{
            \makecell[{{p{\linewidth}}}]{
                \texttt{\tiny{[GM$|$GM]}}
                \texttt{Doubtful: 3} \\
\texttt{Surprised: 0} \\
\texttt{Understood: 2} \\
\texttt{Disgusted: 0} \\
            }
        }
    }
    & & \\ \\

    \theutterance \stepcounter{utterance}  
    & & & \multicolumn{2}{p{0.3\linewidth}}{
        \cellcolor[rgb]{0.9,0.9,0.9}{
            \makecell[{{p{\linewidth}}}]{
                \texttt{\tiny{[GM$|$GM]}}
                \texttt{game\_result = LOSE} \\
            }
        }
    }
    & & \\ \\

    \theutterance \stepcounter{utterance}  
    & & \multicolumn{4}{p{0.6\linewidth}}{
        \cellcolor[rgb]{0.95,0.95,0.95}{
            \makecell[{{p{\linewidth}}}]{
                \texttt{\tiny{[P1$\leftrightarrow$P]}}
                \texttt{Post{-}episode Reflexion (not sent to GM):} \\
\texttt{New plan: I failed to recognize the definitive nature of Jackson's concession. In the dialogue, Jackson explicitly states, "You're right, I've never had to worry about surviving," which is a direct validation of Alejandro's argument about survival. This is not merely a defensive retort or a moment of doubt; it is a clear admission that he understands the other character's perspective. My previous attempts underestimated the intensity of this realization. Next time, I will identify explicit verbal concessions as the strongest indicator of the emotion "Understood." I will assign a score of 10 to "Understood" to reflect this complete shift in perspective, while significantly lowering the scores for "Disgusted" and "Doubtful" to accurately represent that these emotions have been superseded by this moment of clarity.} \\
            }
        }
    }
    & \\ \\

\end{supertabular}
}

\end{document}
